\documentclass[aspectratio=169]{beamer}
\input{../theme/imsb_beamer.tex}

\title{Tissue organisation, communication, and interpretation}
\author{Dr.\ Robin Khatri}
\institute{Institute of Medical Systems Bioinformatics, UKE}
\date{Day 3 \textendash{} 23.06.2026}

\begin{document}

\begin{frame}[plain]\titlepage\end{frame}

\begin{frame}{Today, and why it stands alone}
  \begin{itemize}
    \item We run on a separate, smaller machine (6 GB) and a small, pre-annotated
      single-cell spatial dataset. Nothing today depends on Days 1\textendash 2.
    \item Focus shifts from \emph{making} the data to \emph{reading} it:
      organisation, communication, and careful interpretation.
    \item These are the analyses that turn a cell-type map into biology.
  \end{itemize}
\end{frame}

\begin{frame}{From cell types to tissue organisation}
  \centering
  \begin{tikzpicture}[node distance=10mm]
    \node[box] (a) {cell types\\(who is here)};
    \node[box,right=of a] (b) {neighbourhoods\\(who is next to whom)};
    \node[box,right=of b] (c) {niches\\(recurring units)};
    \node[box,right=of c] (d) {communication\\(signalling)};
    \draw[arr] (a)--(b); \draw[arr] (b)--(c); \draw[arr] (c)--(d);
  \end{tikzpicture}
  \vskip0.6em
  Each step is built from the spatial graph. Proportions alone miss the
  arrangement, which is where tissue function lives.
\end{frame}

\begin{frame}{Neighbourhood enrichment and co-occurrence}
  \begin{itemize}
    \item \textbf{Enrichment}: permute cell-type labels on the graph; compare
      observed adjacency to the null. High means co-location, low means avoidance.
    \item \textbf{Co-occurrence}: how the chance of finding type B near type A
      changes with radius, separating contact from regional association.
  \end{itemize}
  \vskip0.4em
  Caveat: both are relative to the chosen graph and the section's composition, so a
  difference between samples can be a composition difference.
\end{frame}

\begin{frame}{Niches and tissue domains: the method landscape}
  \small
  \begin{itemize}
    \item \textbf{Composition clustering} (what we do): summarise neighbourhood
      cell-type composition, then cluster. Transparent, fast, CPU-friendly.
    \item \textbf{BANKSY}: augments each cell with its neighbourhood; types cells and
      finds domains together, scalable and batch-aware.
    \item \textbf{CellCharter}: clusters neighbourhood-aware embeddings into niches
      across samples.
    \item \textbf{SpaGCN, STAGATE}: graph neural networks for spatial domains.
  \end{itemize}
  Learned methods can capture subtler domains; they add assumptions and compute.
  \refline{Singhal et al., \emph{Nat Genet} 2024 (BANKSY); Varrone et al., \emph{Nat Genet} 2024 (CellCharter); Hu et al., \emph{Nat Methods} 2021 (SpaGCN); Dong \& Zhang, \emph{Nat Commun} 2022 (STAGATE).}
\end{frame}

\begin{frame}{Building niches transparently}
  \centering
  \begin{tikzpicture}[node distance=8mm and 12mm]
    \node[box] (a) {for each cell:\\neighbourhood\\composition vector};
    \node[box,right=of a] (b) {cluster the\\vectors\\(k-means)};
    \node[box,right=of b] (c) {niche labels\\on the tissue};
    \draw[arr] (a)--(b); \draw[arr] (b)--(c);
  \end{tikzpicture}
  \vskip0.6em
  The output is named, mappable units you can compare across samples, with no black
  box between the data and the label.
\end{frame}

\begin{frame}{Cell\textendash cell communication: the logic}
  \begin{columns}[T]
  \begin{column}{0.54\textwidth}
  Score ligand\textendash receptor pairs between cell-type groups: is the ligand
  expressed by one type and the receptor by a co-located type?
  \vskip0.4em
  In spatial data we can add the constraint that the partners are actually
  neighbours, which dissociated data cannot.
  \end{column}
  \begin{column}{0.44\textwidth}
  \centering
  \begin{tikzpicture}[scale=0.95]
    \node[draw=imsbblue,thick,circle,fill=imsbblue!12,minimum size=1.3cm] (s) at (0,0) {sender};
    \node[draw=imsbgrey,thick,circle,fill=imsbgrey!12,minimum size=1.3cm] (r) at (3,0) {receiver};
    \draw[arr] (s) -- node[lab,above]{ligand $\to$ receptor} (r);
    \node[lab,below=1mm of s]{expresses L};
    \node[lab,below=1mm of r]{expresses R};
  \end{tikzpicture}
  \end{column}
  \end{columns}
\end{frame}

\begin{frame}{Communication: methods and caveats}
  \small
  \begin{itemize}
    \item \textbf{Non-spatial:} CellPhoneDB, CellChat, NicheNet (adds downstream
      targets); LIANA unifies many under one interface.
    \item \textbf{Spatial-aware:} squidpy \texttt{ligrec}, stLearn, COMMOT (optimal
      transport), NCEM (niche effects on expression).
  \end{itemize}
  \vskip0.3em
  Caveats: expression is a proxy for activity; a small panel covers few pairs;
  p-values depend on the permutation scheme; segmentation errors move ligands to the
  wrong cell. Read results as hypotheses.
  \refline{Efremova et al., \emph{Nat Protoc} 2020 (CellPhoneDB); Jin et al., \emph{Nat Commun} 2021 (CellChat); Browaeys et al., \emph{Nat Methods} 2020 (NicheNet); Dimitrov et al., \emph{Nat Commun} 2022 (LIANA); Cang \& Nie, \emph{Nat Methods} 2023 (COMMOT); Fischer et al.\ 2023 (NCEM).}
\end{frame}

\begin{frame}{Spatially variable genes and pattern reading}
  \begin{itemize}
    \item Moran's I ranks genes by how smoothly they vary across the tissue.
    \item Top genes often trace anatomy or gradients rather than discrete clusters,
      which complements clustering.
    \item Plot the top hits in space and ask what structure each marks.
  \end{itemize}
\end{frame}

\begin{frame}{Comparing conditions}
  \begin{itemize}
    \item With a disease cohort, the question becomes which niches or compositions
      change between groups.
    \item Composition-aware tests (e.g.\ Milo for differential abundance, scCODA for
      composition) avoid the pitfalls of comparing raw proportions.
    \item Spatial features (niche burden, distances, enrichment scores) can be
      summarised per sample and tested like any covariate.
  \end{itemize}
  \refline{Dann et al., \emph{Nat Biotechnol} 2022 (Milo); B\"uttner et al., \emph{Nat Commun} 2021 (scCODA).}
\end{frame}

\begin{frame}{Reading results critically}
  \begin{itemize}
    \item Segmentation errors propagate into co-expression and communication.
    \item Low cytokine and chemokine detection makes immune states uncertain;
      absence is weak evidence.
    \item Spatial QC metrics are still maturing; thresholds are provisional.
    \item Composition differences between sections can masquerade as enrichment.
      Compare like with like.
  \end{itemize}
\end{frame}

\begin{frame}{The frontier}
  \begin{itemize}
    \item Integrating large H\&E images with the molecular layer, and pathology
      image foundation models for region identity (e.g.\ UNI, Virchow,
      H-optimus), with attention to their failure modes.
    \item Multimodal imaging plus Xenium to give better identity to regions of
      interest.
    \item Agentic and LLM-assisted analysis workflows.
  \end{itemize}
  These are open and moving quickly; today's tools are a foundation, not the ceiling.
\end{frame}

\begin{frame}{Further reading}
  \footnotesize
  \begin{itemize}
    \item Varrone et al. CellCharter. \emph{Nat Genet} 2024; Singhal et al. BANKSY. \emph{Nat Genet} 2024.
    \item Cang \& Nie. COMMOT. \emph{Nat Methods} 2023; Dimitrov et al. LIANA. \emph{Nat Commun} 2022.
    \item Dann et al. Milo. \emph{Nat Biotechnol} 2022.
    \item Squidpy tutorials (neighbourhood enrichment, co-occurrence, ligrec).
  \end{itemize}
\end{frame}

\begin{frame}{Hands-on and discussion: \texttt{day\_3/03\_tutorial.ipynb}}
  Dataset: a small pre-annotated seqFISH section. You will:
  \begin{itemize}
    \item run neighbourhood enrichment and co-occurrence;
    \item identify niches from neighbourhood composition;
    \item score ligand\textendash receptor interactions (offline);
    \item rank spatially variable genes.
  \end{itemize}
  \textbf{Capstone:} pick one finding and read it end to end \textendash{} what it
  claims, whether it is consistent, what could confound it, and what would test it
  next. We discuss as a group.
\end{frame}

\end{document}
